\section{Por qué el Leaderboard Success no equivale a Commonsense}

La mejora de Maieutic todavía ocurre dentro del benchmark, por lo que el ponente vuelve a examinar el benchmark mismo como siguiente paso. Si el conjunto de datos contiene un artifact fijo, el modelo puede aprender «cómo obtener puntaje» sin aprender la capacidad objetivo; esta es la diferencia entre dataset solving y task solving.

\subsection{Una puntuación alta y la fragilidad pueden coexistir}

La slide coloca los resultados de alta puntuación en object recognition, question answering y captioning junto a los fallos adversarial u out-of-domain. La comparación horizontal no busca negar el avance, sino señalar que la distribución de evaluación suele ser más estrecha que el uso real. La conclusión inferior, «solo se resuelve el dataset, no la underlying task», exige que los investigadores reporten perturbation, domain shift y failure category, y no solo el leaderboard rank.

\lecturefigure{slide-09-brittle-leaderboards-dataset-task.jpg}{Superhuman leaderboard performance y adversarial brittleness coexisten}{Diapositiva V009 recuperada de la grabación oficial; Stanford CS25 Lecture 16}

\teachervoice{Aquí se concreta el «choose your battles» del ponente: si el benchmark no puede distinguir entre surface pattern y conceptual understanding, seguir acumulando puntaje solo significa ganar en el campo de batalla equivocado. El diseño de nuevas task forma parte de la investigación algorítmica.}

\subsection{Lo que falta entre Human y AI no son más oraciones}

El diagrama de puente coloca a la izquierda la vida cotidiana, el lenguaje y la experiencia social de las personas, y a la derecha las salidas variadas pero inestables de la máquina. El gap del centro apunta al practical knowledge: saber que la puerta del refrigerador normalmente debe estar cerrada, saber por qué una frase resulta ofensiva en una relación determinada, saber qué consecuencias físicas puede tener una acción. Lo que el modelo necesita no son solo más text tokens, sino conocimiento situacional componible y verificable.

\lecturefigure{slide-10-commonsense-gap.jpg}{Commonsense gap: del lenguaje superficial a la comprensión del mundo cotidiano}{Diapositiva V010 recuperada de la grabación oficial; Stanford CS25 Lecture 16}

\subsection{La Operational Definition adoptada en esta clase}

La slide de definición restringe el commonsense a practical knowledge and reasoning, centrado en everyday situations and events, y normalmente compartido por la mayoría de las personas. Las tres restricciones importan: no es la totalidad del conocimiento enciclopédico, no es solo un conjunto de fact y tampoco es absolutamente uniforme entre culturas. Al leer la figura, «commonly shared» debe entenderse como un alcance empírico, no como una verdad moral ni una universal law.

\lecturefigure{slide-11-definition-of-commonsense.jpg}{La definición de sentido común de esta clase: conocimiento y razonamiento cotidianos, prácticos y ampliamente compartidos}{Diapositiva V011 recuperada de la grabación oficial; Stanford CS25 Lecture 16}

La siguiente slide distingue además su utilidad para las personas y para las máquinas. Para las personas, el sentido común sostiene la convivencia y la seguridad; para la AI, el sentido común ayuda a comprender las necesidades del usuario, las consecuencias de las acciones y las condiciones implícitas. El ejemplo de la puerta del refrigerador muestra que la misma frase, «mantener la puerta abierta», puede ser razonable o peligrosa según la situación; por ello, el sentido común no es un phrase lookup fijo, sino una context-conditioned inference.

\lecturefigure{slide-12-human-ai-commonsense-importance.jpg}{Por qué el sentido común importa tanto para la coordinación humana como para que la AI comprenda las acciones}{Diapositiva V012 recuperada de la grabación oficial; Stanford CS25 Lecture 16}

\begin{importantbox}{Tres niveles: Dataset, Task y Capability}
El dataset es una muestra finita; la task es el mapeo que se espera que el sistema realice; la capability es una capacidad estable que se reutiliza entre distribuciones. Una puntuación alta en el conjunto de prueba, como máximo, demuestra directamente dataset-level performance. Para afirmar una commonsense capability se necesitan además reformulaciones diversas, variaciones composicionales, contrafácticos, pruebas entre dominios y failure analysis.
\end{importantbox}

\subsection{Resumen de la sección}

El primer problema de la investigación sobre sentido común es definir el objeto de evaluación. Las tablas de clasificación no son inútiles, pero deben acompañarse de adversarial y out-of-domain evidence. El valor de ATOMIC, COMET y Delphi, que se presentan a continuación, consiste precisamente en convertir «el modelo quizá lo sabe implícitamente» en datos, relaciones e interfaces de tarea más explícitos.
